\chapter{The Retro Ladder: DirectX 5--8 and Glide}
\label{ch:retro-ladder}

The second retro chapter follows four successive Direct3D interface generations and one
contemporary rival.  Every family maps CNA's modern contract onto fixed-function hardware
concepts; none has a programmable shader stage in its declared scope.

\begin{longtable}{p{0.15\textwidth}p{0.28\textwidth}p{0.24\textwidth}p{0.19\textwidth}}
\toprule
Identity & API shape & Extended draws & RT / MRT / query \\
\midrule
\endfirsthead
\toprule
Identity & API shape & Extended draws & RT / MRT / query \\
\midrule
\endhead
\texttt{DIRECTX5} & DirectDraw~4 + Direct3D~3 & native fixed function & yes / no / no \\
\texttt{DIRECTX6} & same COM generation, real stencil & native fixed function & yes / no / no \\
\texttt{DIRECTX7} & flattened Direct3D~7 device model & native fixed function & yes / no / no \\
\texttt{DIRECTX8} & Direct3D~8 device API & native fixed function & yes / no / no \\
\texttt{GLIDE} & Glide~3 runtime ABI & native fixed function & no / no / no \\
\bottomrule
\end{longtable}

There is no \texttt{DIRECTX4} selector because DirectX~4 was not a public release.  The jump
in names is historical, not a missing CNA family.

\section{DIRECTX5: the stable immediate-mode baseline}

\texttt{DIRECTX5} moves to DirectDraw~4 and Direct3D~3, where flexible vertex formats and
\texttt{DrawPrimitive} form the baseline.  Execute buffers are fully absent from CNA's path.
The earlier spike result therefore becomes a repeatable port strategy rather than a one-off
experiment.

The renderer translates effect-aware draws into fixed-function matrices, texture stages,
lighting, fog, alpha test, blend, depth, and raster state.  A native override does not mean
arbitrary effects are possible; it means supported stock semantics are reduced intentionally
inside the family instead of being dropped by the interface default.  2D targets work, while
MRT and occlusion queries remain outside the era's contract.

\section{DIRECTX6: a small delta with a real stencil gain}

The DirectX~6 port retained the same broad COM interface shape as DirectX~5.  The measured
source delta was about six percent, an example of why close lineage is useful: shared behavior
can be brought up first, then the changed surface can be audited directly.

One change is user-visible.  This rung implements and tests real stencil behavior.  Its
\texttt{SupportsCapability(StencilBuffer)} answer nevertheless returns false.  That is a
capability-reporting defect, not evidence that the implementation lacks stencil.  Callers
that trust the bit will unnecessarily disable a working feature; callers that ignore it may
use behavior the renderer demonstrably provides.

\section{DIRECTX7: the flattened device model}

Direct3D~7 removes the separate viewport-object model and folds more state into the device.
The CNA contract above it remains stable, which makes this renderer a useful adapter study:
the same viewport and state calls require a different ownership mapping even when the visible
scene is unchanged.

Its feature outline remains native fixed-function draws, 2D targets, no MRT, and no occlusion
queries.  Real stencil exists, while the same false capability answer found in DirectX~6
persists here.

\section{DIRECTX8: a new device generation, deliberately no Shader Model 1}

Direct3D~8 is a much more recognizable device API, but CNA intentionally stops at its
fixed-function surface.  Shader Model~1 assembly is too narrow for realistic XNA effect
content and would create a misleading partial programmable promise.  The renderer therefore
keeps the same honest semantic ceiling as the preceding rungs.

The build has a distinctive dependency.  A normal 64-bit mingw-w64 installation does not
provide the required D3D8 import library; CNA expects the D8VK/DXVK import path and refuses
configuration when it is absent.  Runtime evidence uses a dedicated Wine prefix and can
force llvmpipe through the DXVK device filter for reproducibility.

2D targets are implemented; MRT and queries are not.  Stencil is implemented but its
capability answer is false, completing the same defect across DirectX~6, 7, and 8.

\begin{sourcenote}
For all four DirectX identities, a passing Wine run proves CNA plus Wine or DXVK.  It does not
retroactively prove every native driver shipped with the historical SDK.  The compatibility
layer is part of the recorded environment.
\end{sourcenote}

\section{GLIDE: a contemporary fixed-function branch}

\texttt{GLIDE} targets the Glide~3 ABI and is intentionally constrained to Windows and
32-bit builds.  CNA does not vendor the runtime DLL; a user-supplied \texttt{glide3x.dll} is
required.  That packaging boundary exists partly because common compatibility implementations
have licenses unsuitable for vendoring here.

The renderer owns a native fixed-function extended draw path, but it cannot create
\cnaclass{RenderTarget2D}; the factory throws.  MRT and occlusion queries are likewise absent.
This makes Glide the only family in the retro table without even the single off-screen target
that the DirectX rungs provide.

Its evidence boundary is sharper still.  The repository contains a real example target, but
deliberately does not register it with \texttt{ctest}; no automated job executes Glide
rendering.  Compilation can validate headers and ABI shape, while pixels remain a manual
claim requiring a suitable 32-bit runtime.

\section{What the ladder teaches}

The sequence is less about nostalgia than contract design:

\begin{itemize}
  \item a single public draw packet can be reduced across several fixed pipelines without
        pretending that arbitrary shaders exist;
  \item measured port similarity accelerates implementation but cannot replace per-family
        capability and presentation evidence;
  \item capability bits can be less accurate than executable behavior, as the three stencil
        false negatives demonstrate;
  \item build, compatibility-layer, and runtime-DLL success are distinct proof tiers.
\end{itemize}

\section{A practical retro verification pass}

For each DirectX rung, exercise clear/present, one textured 2D draw, one depth-ordered 3D
scene, target bind/unbind, and every state feature the family claims.  On DirectX~6--8 add a
stencil pixel test and record the known capability mismatch.  Capture the Wine version,
translation layer, prefix, and device name beside the result.

For Glide, keep compile evidence and execution evidence separate.  A successful 32-bit link
does not show that a compatible DLL loaded, and a loaded DLL does not show that a triangle
reached the screen.  Until the example is executed and its pixels observed in a documented
environment, the honest status remains implemented ABI with unautomated rendering evidence.
